\section{Conclusion}\label{sec:conclusion}

We studied how a UAV swarm without infrastructure can get victim reports to responders in time.
SENTINEL lets UAVs search freely, triggers delivery from the information each UAV has confirmed,
fuses views so that occluded angles cannot veto a clear detection, and preserves the links of a
delivery path for every feasible step (Proposition~\ref{prop:link}). In about 13\,000 simulated missions against seven
tuned baselines, it reported significantly more victims within 60\,s than every baseline in most
settings (urban, suburban and open channels, all failure settings in urban rubble, swarms of 3--50
UAVs on the 400\,m area and moving obstacles), it stayed significantly ahead of the two baselines
tested under detector uncertainty, in the severe channel it was ahead of or tied with the best
baseline, and it was never significantly worse than any of the seven baselines. A PX4/Gazebo cross-check
agreed with the fast simulator within 0.07.
The ablation found significant effects for the fusion rule and the delivery trigger; in the urban
400\,m scenarios SENTINEL remained significantly ahead when the baselines were given its fusion rule,
whereas in a dense 50-UAV swarm the fusion rule alone explained its advantage. Continuous connectivity
did not help, and the link-preservation layer reduced preserved-link violations 37-fold. Future work will adapt the delivery trigger to the observed contact rate,
calibrate the channel model with measurements, and validate the system in hardware flights.
